% ---------------------------------------------------------------------------
% HPOP — CPA annotation: theory + baseline evaluation
% Self-contained; compile with pdflatex. Drop the two \section blocks into the paper.
% ---------------------------------------------------------------------------
\documentclass[11pt]{article}
\usepackage{amsmath,amssymb}
\usepackage{booktabs}
\usepackage{array}
\usepackage[margin=1in]{geometry}
\newcommand{\cpa}{\textsc{cpa}}
\begin{document}

\section{Procedural annotation framework}
\label{sec:framework}

We annotate each agent trajectory in three layers, following a standoff multi-layer
design conformant to the Linguistic Annotation Framework
(ISO~24612{:}2012; \cite{bird1999}).

\paragraph{Layer B$\cdot$1a: action tokens.}
Each raw event is normalised, with \emph{no} interpretive rule, into an
\emph{action token} carrying four fields,
$a = (\textit{action\_type},\,\textit{tool\_family},\,\textit{agent\_id},\,\textit{artifact\_id})$,
and adjacent identical tokens are merged (de-duplication). This is pure data cleaning
and is the substrate every downstream layer references.

\paragraph{Layer A: cognitive phases.}
Each span is assigned, \emph{first}, one of nine cognitive control states grounded in
the prefrontal-cortex theory of cognitive control \cite{miller2001}. The phases are
organised over four abstraction levels --- L1 goal-setting (\textsc{plan}); L2 execution
(\textsc{retrieve}, \textsc{inspect}, \textsc{extract}, \textsc{verify}, \textsc{write});
L3 integration (\textsc{synthesize}); L4 interrupt (\textsc{repair}, \textsc{handoff}).
The phase label is an \emph{optional, coarse procedural descriptor} used for interpretability and
diagnostic analysis; it is \emph{not} a fixed skill type. HPOP does not take phases as skill labels:
it receives the \cpa{} occurrence sequences and learns a reusable library of \emph{local
partial-order skills}, and a learned skill may contain \cpa s from one phase or several. The phase
taxonomy is given in Table~\ref{tab:phases}.

\begin{table}[t]
\centering\small
\begin{tabular}{@{}l l p{5.4cm} l r@{}}
\toprule
Level & Phase & Cognitive control state & Grounding & occ. \\
\midrule
L1 goal-setting & \textsc{plan}       & goal decomposition; PFC sets the goal hierarchy            & \cite{miller2001,grosz1986}      & 0 \\
\addlinespace[1pt]
L2 execution    & \textsc{retrieve}   & directed, externally-directed acquisition                 & \cite{botvinick2009}             & 4{,}298 \\
L2 execution    & \textsc{inspect}    & metacognitive monitoring; no proposition tested           & \cite{nelson1990}                & 4{,}079 \\
L2 execution    & \textsc{extract}    & central executive locks on a target value                 & \cite{baddeley2000}              & 1{,}681 \\
L2 execution    & \textsc{verify}     & metacognitive control; output is a truth value            & \cite{nelson1990,wason1960}      & 7{,}931 \\
L2 execution    & \textsc{write}      & generative production; externalisation mode               & \cite{flower1981}                & 5{,}824 \\
\addlinespace[1pt]
L3 integration  & \textsc{synthesize} & multi-source integration into a new representation        & \cite{hutchins1995}              & 0 \\
\addlinespace[1pt]
L4 interrupt    & \textsc{repair}     & error monitoring; SAS override after a failure signal     & \cite{norman1986}                & 116 \\
L4 interrupt    & \textsc{handoff}    & distributed cognitive load transfer (agent change)        & \cite{hutchins1995,monsell2003}  & 0 \\
\bottomrule
\end{tabular}
\caption{The nine cognitive phases over four abstraction levels, with their control-state
definitions and grounding. The \emph{occ.} column gives the per-phase occurrence count over the
500-trajectory corpus. The empty L3/L4 phases are a finding, not a gap: single-agent coding traces
do not integrate across sources (\textsc{synthesize}) or transfer control (\textsc{handoff}), and
overt planning (\textsc{plan}) is latent in reasoning the trace does not externalise.}
\label{tab:phases}
\end{table}

\paragraph{Layer B$\cdot$1b: canonical procedural actions.}
Within each phase, the span is refined to a \emph{canonical procedural action} (\cpa) ---
the procedural \emph{function} realising that control state, more abstract than a tool call
and portable across repositories. \cpa s correspond to the \emph{Action} tier of Activity
Theory \cite{leontiev1978}: goal-directed, consciously executed units, as opposed to the
automatised \emph{Operations} (raw clicks/commands) at which reusable structure does not emerge.
The annotation is thus coarse-to-fine: phase first, \cpa\ within it, each \cpa\ realising
exactly one phase.

\section{Evaluating the annotation against no-instruction baselines}
\label{sec:eval}

A labelling scheme assigns a symbol $\ell\in V$ to each occurrence. A \emph{good} scheme makes
the trajectory predictable: the current label should be informative about the next. We
therefore measure the sequential structure a scheme induces, relative to a null that
destroys order.

\subsection{Transition information and the shuffle null}
For a scheme, collect the multiset of consecutive label pairs over all trajectories,
$\mathcal{P}=\{(\ell_t,\ell_{t+1})\}$, and estimate the \emph{transition mutual information}
\begin{equation}
\hat{I} \;=\; \sum_{x,y\in V} \hat{p}(x,y)\,
   \log_2\!\frac{\hat{p}(x,y)}{\hat{p}(x)\,\hat{p}(y)} ,
\label{eq:mi}
\end{equation}
where $\hat p(x,y)$ is the empirical joint over adjacent pairs and $\hat p(x),\hat p(y)$ the
current/next marginals. If the next label is independent of the current one, $\hat I=0$.

Finite samples inflate $\hat I$ by an $O(|V|^2/N)$ bias, so a larger vocabulary scores
spuriously higher. We remove it with a within-trajectory \emph{shuffle null}: permute each
trajectory's labels uniformly at random (preserving length and unigram composition, destroying
order), recompute \eqref{eq:mi}, and average over $B=150$ permutations to obtain $\bar{I}_0$.
Because the null shares the exact marginals and $N$ of the data, $\bar I_0$ estimates the
finite-sample floor for that scheme. We report the \emph{excess} structure and a
size-normalised \emph{efficiency},
\begin{equation}
\Delta \;=\; \hat{I}-\bar{I}_0 ,
\qquad
\eta \;=\; \frac{\Delta}{\log_2 |V|},
\label{eq:excess}
\end{equation}
where $\log_2|V|$ is the maximum information a single symbol can carry, making schemes of
different vocabulary sizes comparable. By construction a labelling with no ordering structure
has $\Delta\approx 0$.

\subsection{Baselines}
All schemes label the \emph{same} occurrence segmentation; only the symbol differs.
\begin{description}
  \item[\textsc{cmdword} (no instruction).] $\ell$ is the raw leading command token
        (\texttt{pytest}, \texttt{grep}, \texttt{view}, \dots); no taxonomy. A natural,
        uncontrolled vocabulary.
  \item[\textsc{actobj} (no instruction).] $\ell=(\textit{action\_type},\textit{artifact\_role})$,
        e.g.\ \texttt{run:test}, \texttt{edit:source}; still no procedural taxonomy.
  \item[v1 \cpa\ (induced).] The bottom-up dictionary: surface forms mined from data, then
        grouped into phases post hoc.
  \item[v2 \cpa\ (cognitive).] The top-down dictionary of Section~\ref{sec:framework}: nine
        phases first, \cpa s derived as their realisations.
\end{description}

\subsection{Results}
Table~\ref{tab:ladder} reports the ladder on 500 SWE-rebench/OpenHands trajectories
(334 repositories; $23{,}929$ occurrences). A \emph{no-instruction} labelling is the honest
floor; both principled dictionaries lie well above it, and the cognitive v2 scheme attains the
highest excess and the highest per-symbol efficiency.

\begin{table}[t]
\centering
\begin{tabular}{l c c c c c}
\toprule
scheme & $|V|$ & $\hat I$ & $\bar I_0$ & $\Delta$ (bits) & $\eta$ \\
\midrule
\textsc{cmdword} (no instruction) & 62 & 0.623 & 0.041 & 0.582 & 0.098 \\
\textsc{actobj}\ (no instruction) & 24 & 0.941 & 0.023 & 0.918 & 0.200 \\
v1 \cpa\ (induced)                & 28 & 0.968 & 0.027 & 0.940 & 0.196 \\
\textbf{v2 \cpa\ (cognitive)}     & 29 & 1.009 & 0.030 & \textbf{0.979} & \textbf{0.201} \\
\bottomrule
\end{tabular}
\caption{Sequential structure beyond a shuffled null (Eq.~\ref{eq:excess}). The naive
no-instruction labelling sprawls to $|V|=62$ yet is half as efficient ($\eta=0.098$); the
cognitive dictionary packs the most non-random structure into the fewest reusable symbols.}
\label{tab:ladder}
\end{table}

\paragraph{Information beyond the surface.}
To confirm the abstraction is not a relabelling of the physical layer, we compare three
\emph{layers} of the same trace by (i) normalised predictability
$\mathrm{NP}=\hat I(\ell_t;\ell_{t+1})/\hat H(\ell_t)$ and (ii) the information a current label
gives about the next physical action,
$\mathrm{IG}=\hat H(a_{t+1})-\hat H(a_{t+1}\mid \ell_t)$. At nearly equal vocabulary the
cognitive phase layer ($|V|{=}6$, $\mathrm{NP}{=}0.186$, $\mathrm{IG}{=}0.253$\,b) is about
twice as self-predictable as the raw physical layer ($|V|{=}5$, $\mathrm{NP}{=}0.097$,
$\mathrm{IG}{=}0.159$\,b), and a current phase is \emph{more} predictive of the next physical
action than the current physical action itself --- the abstraction explains the surface,
not vice versa.

\begin{thebibliography}{9}
\bibitem{bird1999} S.~Bird and M.~Liberman. A formal framework for linguistic annotation.
  \emph{Speech Communication}, 1999.
\bibitem{miller2001} E.~K.~Miller and J.~D.~Cohen. An integrative theory of prefrontal cortex
  function. \emph{Annual Review of Neuroscience}, 2001.
\bibitem{leontiev1978} A.~N.~Leontiev. \emph{Activity, Consciousness, and Personality}. 1978.
\bibitem{grosz1986} B.~J.~Grosz and C.~L.~Sidner. Attention, intentions, and the structure of
  discourse. \emph{Computational Linguistics}, 1986.
\bibitem{botvinick2009} M.~M.~Botvinick, Y.~Niv, and A.~C.~Barto. Hierarchically organized
  behavior and its neural foundations. \emph{Cognition}, 2009.
\bibitem{nelson1990} T.~O.~Nelson and L.~Narens. Metamemory: A theoretical framework and new
  findings. \emph{Psychology of Learning and Motivation}, 1990.
\bibitem{baddeley2000} A.~Baddeley. The episodic buffer: a new component of working memory?
  \emph{Trends in Cognitive Sciences}, 2000.
\bibitem{wason1960} P.~C.~Wason. On the failure to eliminate hypotheses in a conceptual task.
  \emph{Quarterly Journal of Experimental Psychology}, 1960.
\bibitem{flower1981} L.~Flower and J.~R.~Hayes. A cognitive process theory of writing.
  \emph{College Composition and Communication}, 1981.
\bibitem{hutchins1995} E.~Hutchins. \emph{Cognition in the Wild}. MIT Press, 1995.
\bibitem{norman1986} D.~A.~Norman and T.~Shallice. Attention to action: Willed and automatic
  control of behavior. In \emph{Consciousness and Self-Regulation}, 1986.
\bibitem{monsell2003} S.~Monsell. Task switching. \emph{Trends in Cognitive Sciences}, 2003.
\end{thebibliography}

\end{document}
